\subsection{Searching for Extra Dimensions and Copies of the Standard Model with IceCube --- arXiv:2608.29746}

\textbf{Authors:} IceCube Collaboration

\paragraph{主な主張}
IceCube の 10.7 年分の上向きミューニュートリノ事象を用いて large extra dimension および Standard Model copies を探索し、最大の余剰次元半径に対して $R\lesssim 0.17\,\mu\mathrm{m}$（90\% C.L.）という制限を得た \cite{IceCube:2026extra}。

\paragraph{新規性}
TeV スケールの大気ニュートリノが地球内部を通過する際の matter-enhanced active--sterile oscillation を利用して KK neutrino tower に感度を持たせ、parameter space の一部で従来最強または未探索領域への制限を与えた点が新しい。

\paragraph{位置付け}
重力実験ではなくニュートリノ振動から sub-micron compactification scale を制限する phenomenological probe であり、bulk right-handed neutrino を含む large-extra-dimension realization に対する直接的な実験制約として位置付けられる。

\paragraph{コメント：$T^2$ への読み替え}
矩形 torus では最軽 KK mode は概ね最大半径 $R_{\max}$ に対して $m_{\mathrm{KK,min}}\sim 1/R_{\max}$ で決まるが、一般の $T^2$ では $m_{\mathrm{KK}}^2=n_a G^{ab}(T,\tau)n_b$ なので、IceCube の制限は本質的には $m_{\mathrm{KK,min}}^2=\min_{\mathbf n\neq0}n_aG^{ab}(T,\tau)n_b$ への下限として、K\"ahler modulus と complex structure modulus の許容領域へ写像するのが自然である。したがって $R<0.17\,\mu\mathrm{m}$ をそのまま $\tau$ の制限と読むのではなく、bulk 右巻きニュートリノが同じ compact space を伝播する具体模型を定めた上で KK 固有値スペクトルを通じて $(T,\tau)$ 空間の exclusion region を構成する必要がある。
